\section{Introduction}\label{sec:introduction}

% Target: ~1.5 two-column pages (~3 single-column equivalent)

% [Paragraph 1: Problem statement and motivation.]
%   - Define information channels; introduce timing channels as a covert
%     communication mechanism that bypasses OS security enforcement.
%   - Cite real-world attacks: Dhem et al. (2000), Aciicmez (2005) for crypto
%     side-channels; Ristenpart et al. (2009) for cross-VM leaks in cloud
%     environments; Lipp et al. (2018) for Meltdown; Kocher et al. (2019) for
%     Spectre and speculative-execution attacks.
%   - Emphasise that the OS scheduler itself can act as a timing channel by
%     making task execution time observable across security domains.

% [Paragraph 2: Existing mitigations and their limitations.]
%   - User-level approaches: constant-time cryptography (FaCT, Cauligi et al.
%     2019), kernel page-table isolation (Corbet 2017). Burden on developer;
%     cannot address scheduling-level or microarchitectural leaks.
%   - Clean-slate secure microkernels: seL4 time protection (Ge et al. 2019),
%     S3K (Karlsson et al. 2025). Built from ground up with formal guarantees;
%     not deployable on existing mainstream OS without a complete rewrite.
%   - Formal verification gap: seL4 memory isolation is formally verified
%     (Murray et al. 2013), but no kernel is formally verified against timing
%     channels (Buckley et al. 2023).

% [Paragraph 3: Research gap and research question.]
%   - Ge et al. hint that time protection could be implemented in other systems,
%     but retrofitting into an existing kernel is fundamentally different from
%     designing from scratch.
%   - FreeRTOS as target: widely deployed, priority-based preemptive scheduler
%     inherently vulnerable to scheduling-based timing channels.
%   - Research question: \emph{Can time protection mechanisms be integrated into
%     a mainstream OS to mitigate timing channels?}

% [Paragraph 4: Approach and contributions.]
%   - Defence mechanism: domain-level scheduler that partitions execution into
%     fixed time slices, combined with deterministic microarchitectural state
%     flushing via the RISC-V temporal fence (fence.t) instruction.
%   - Three contributions:
%     \begin{itemize}
%       \item A domain-level scheduler built on FreeRTOS-MPU that decouples
%         task schedulability across security domains.
%       \item Temporal isolation through fence.t flushing on every domain
%         switch, preventing shared microarchitectural state leakage.
%       \item An evaluation on QEMU system emulation with instruction-based
%         cycle counting, demonstrating constant-time domain scheduling with
%         negligible variance.
%     \end{itemize}

% [Paragraph 5: Paper outline.]
%   - Section II: Background on timing channels, isolation, fence.t, and
%     FreeRTOS.
%   - Section III: Related work (user-level mitigations, seL4, S3K).
%   - Section IV: Design and analysis of the domain-level scheduler.
%   - Section V: Implementation details.
%   - Section VI: Evaluation results.
%   - Section VII: Discussion and future work.
%   - Section VIII: Conclusion.
